\documentclass[11pt]{article}
\usepackage[margin=1in]{geometry}
\usepackage{amsmath,amssymb,amsthm,booktabs,hyperref,microtype}
\newtheorem{proposition}{Proposition}
\newtheorem{theorem}{Theorem}
\title{Lambda-Square Shadows of a \texorpdfstring{$W(E_6)$}{W(E6)} Gassmann Pair}
\author{C62 computational research note}
\date{August 2026}
\begin{document}
\maketitle

\begin{abstract}
Starting from the released degree-160 Gassmann pair for the Weyl group
$G=W(E_6)$, we compute the exterior-square and symmetric-square shadows of
the two permutation $G$-sets.  Exact finite-group enumeration verifies equal
rational characters, complete orbit atlases, core-free stabilizers, normalizer
orders, and a fixed-field dictionary grouped by ambient conjugacy.  The two
shadows have dimensions $51{,}040$ and $51{,}360$.  Product-form marker
carriers are supplied with a split-prime noncollision witness.  These are
finite-group and formal marker results only: no arithmetic field resolvent,
discriminant, Euler factor, or root-number claim is made.
\end{abstract}

\section{Question and scope}
Let $H_+$ and $H_-$ be the released nonconjugate order-$162$ subgroups of
$G=W(E_6)$, and write $X_\pm=G/H_\pm$.  The question is whether applying the
degree-two lambda operations can preserve the permutation character while
separating the resulting finite $G$-sets and their fixed-field shadows.
Throughout, the scope literal is
\texttt{NO\_BAD\_EULER\_OR\_ROOT\_NUMBER}.  Thus all conclusions below are
finite-group or formal-marker conclusions.

\section{Lambda identities}
For a finite $G$-set $X$ and $g\in G$, the standard formulas are
\[
 \chi_{\Lambda^2 X}(g)=\frac{\chi_X(g)^2-\chi_X(g^2)}2,
 \qquad
 \chi_{\operatorname{Sym}^2X}(g)=\frac{\chi_X(g)^2+\chi_X(g^2)}2.
\]
The released C61 equality $\chi_{X_+}=\chi_{X_-}$ therefore implies equality
of both degree-two characters.  The producer independently evaluates these
identities on the complete $W(E_6)$ action and records dimensions
\[
 |\Lambda^2X_\pm|=\binom{320}{2}=51{,}040,
 \qquad |\operatorname{Sym}^2X_\pm|=\binom{321}{2}=51{,}360.
\]

\section{Complete orbit atlases}
The exact action enumerates all two-subsets and all size-two multisets.  The
exterior-square atlas has 10 orbits totaling $51{,}040$ points; the
symmetric-square atlas has 11 orbits totaling $51{,}360$ points.  For every
orbit we store the complete stabilizer element set, its core, its normalizer,
and canonical digests.  In each row the fixed-field degree is checked by
\[
 [G^{\mathrm{core}(S)}:S]=\frac{|G|}{|S|}=\frac{51{,}840}{|S|},
\]
with trivial core in all 42 plus/minus orbit records.

Four matched exterior rows and five matched symmetric rows have
nonconjugate plus/minus stabilizers.  Hence character equality is not used as
an isomorphism test for the actual finite $G$-sets.

\section{Fixed-field dictionary}
The G4 dictionary groups the complete stabilizer sets by conjugacy in the
ambient $W(E_6)$, producing 16 explicit type labels.  The checker verifies
that each label has a consistent subgroup order, normalizer order, core order,
and fixed-field degree.  The plus and minus type sets differ in both lambda
operations.  Several subgroup orders split into multiple labels, so subgroup
order, orbit-table position, or a hash cannot serve as the field identifier.

\begin{center}
\begin{tabular}{lrrr}
\toprule
shadow & orbit count & total degree & matched nonconjugate rows\\
\midrule
$\Lambda^2$ & 10 & 51,040 & 4\\
$\operatorname{Sym}^2$ & 11 & 51,360 & 5\\
\bottomrule
\end{tabular}
\end{center}

\section{Marker carriers and arithmetic boundary}
For each orbit $\mathcal O$ the reproducibility package stores the formal
carrier
\[
 R_{\mathcal O}(T)=\prod_{(i,j)\in\mathcal O}
 \bigl(T-(512X_i+X_j)\bigr).
\]
At the released split-prime witness $p=692717$, all evaluated marker values
inside each orbit are distinct and below $p$.  This certifies a convenient
noncollision label for the orbit computation.  It does not expand the carrier
over characteristic zero and does not identify an arithmetic number field.
Consequently no discriminant, different, local extension, Euler factor,
ramification, or root-number statement is made here.

\section{Reproducibility and nonclaims}
The code and JSON evidence are deterministic and source-bound to the released
C61 group action.  The commands are:
\begin{verbatim}
python3 code/c62_lambda.py
python3 code/c62_atlas.py
python3 code/c62_resolvent.py
python3 code/c62_dictionary.py
python3 code/c62_dictionary_checker.py
\end{verbatim}
All outputs are marked prefreeze until the final audit and manifest closure.
The paper deliberately does not claim arithmetic equivalence beyond the
finite fixed-field shadow, and it does not claim any bad-prime or root-number
result.

\section{Conclusion}
Degree-two lambda operations preserve the rational permutation character of
the C61 Gassmann pair while producing complete, explicitly separated finite
orbit/stabilizer data.  The resulting 16-type fixed-field dictionary and the
formal marker carriers constitute the C62 contribution within the stated
scope boundary.

\bibliographystyle{plain}
\begin{thebibliography}{9}
\bibitem{serre}
J.-P. Serre, \emph{Topics in Galois Theory}, Jones and Bartlett, 1992.
\bibitem{burnside}
W. Burnside, \emph{Theory of Groups of Finite Order}, Cambridge University Press, 1911.
\end{thebibliography}
\end{document}
